\chapter*{Résumé}
\addcontentsline{toc}{chapter}{Résumé}
\thispagestyle{plain}

Ce rapport présente la conception et la réalisation d'un \textbf{jumeau numérique d'une ligne de production simplifiée}, dans le cadre du mini-projet Industrie~4.0 de l'\ac{ENSAM} Rabat. La ligne étudiée comporte quatre postes successifs --- alimentation (5~s), traitement (8~s), contrôle qualité (15~s) et emballage (7~s) --- reliés par un convoyeur, avec des files d'attente de capacité limitée, des pannes aléatoires et un taux de rejet de 5~\% au contrôle.

Le système repose sur une architecture en trois briques. Un \textbf{modèle de simulation} écrit en Python reproduit le comportement de la ligne par pas de temps de 0{,}05~s~: états des machines (en marche, en attente, bloquée, en panne), blocages de l'amont, propagation des pannes et calcul des \ac{KPI}. Un \textbf{backend} FastAPI diffuse l'état complet de la ligne dix fois par seconde par \ac{WS}, reçoit les commandes de l'opérateur et enregistre l'historique de production dans une base SQLite. Une \textbf{interface} web affiche la ligne en \ac{3D} avec une scène Unity exportée en WebGL et un tableau de bord de supervision synchronisé sur le même flux de données.

Le scénario d'optimisation porte sur le goulot d'étranglement du contrôle qualité. Sur une heure simulée avec pannes et une graine aléatoire identique pour tous les scénarios, la ligne initiale produit \textbf{218 produits}. Un contrôle automatisé à 8~s ou un second poste de contrôle en parallèle portent la production à \textbf{393} et \textbf{394 produits} (+80~\%), divisent l'attente au contrôle par 9 à 13 et déplacent le goulot vers le poste de traitement. Ces résultats sont mesurés avec la scène Unity connectée au backend. Un convoyeur adaptatif ne change pas le débit mais réduit l'en-cours moyen de 51~\%. La cohérence du modèle est vérifiée par 23 tests automatisés (débit de régime égal au débit théorique à 5~\% près, conservation des produits, capacité des files, effet des pannes). Les limites --- données entièrement simulées, absence de connexion à des équipements réels --- sont discutées en conclusion.

\vspace{0.4cm}
\textbf{Mots-clés~:} Industrie~4.0, jumeau numérique, simulation de ligne de production, goulot d'étranglement, supervision temps réel, WebSocket, Unity, WebGL, indicateurs de performance.

\vspace{1.2cm}

\chapter*{Abstract}
\addcontentsline{toc}{chapter}{Abstract}
\thispagestyle{plain}

This report presents the design and implementation of a \textbf{digital twin of a simplified production line}, carried out as part of the Industry~4.0 mini-project at \ac{ENSAM} Rabat. The line has four successive stations --- feeding (5~s), processing (8~s), quality control (15~s) and packaging (7~s) --- linked by a conveyor, with limited-capacity buffers, random breakdowns and a 5\% reject rate at quality control.

The system is built from three components. A \textbf{simulation model} written in Python reproduces the behaviour of the line in 0.05~s time steps: machine states (running, idle, blocked, down), upstream blocking, failure propagation and \ac{KPI} computation. A FastAPI \textbf{backend} streams the full line state ten times per second over \ac{WS}, receives operator commands and stores the production history in an SQLite database. A web \textbf{interface} renders the line in \ac{3D} with a Unity scene exported to WebGL, together with a supervision dashboard synchronised on the same data stream.

The optimisation scenario targets the quality-control bottleneck. Over one simulated hour with breakdowns and the same random seed for every scenario, the initial line produces \textbf{218 products}. An automated 8~s inspection or a second parallel inspection station raise output to \textbf{393} and \textbf{394 products} (+80\%), cut the waiting time at inspection by a factor of 9 to 13, and shift the bottleneck to the processing station. These results were measured with the Unity scene connected to the backend. An adaptive conveyor leaves throughput unchanged but reduces average work-in-progress by 51\%. Model consistency is checked by 23 automated tests (steady-state throughput within 5\% of theory, product conservation, buffer capacity, failure effects). The limitations --- fully simulated data, no connection to real equipment --- are discussed in the conclusion.

\vspace{0.4cm}
\textbf{Keywords:} Industry~4.0, digital twin, production line simulation, bottleneck, real-time supervision, WebSocket, Unity, WebGL, performance indicators.
